\begin{table}[H]
\floatleft
\caption{Significant Domain-conditional Trait Moderators}
\label{tab:bydomain}
\vspace{2pt}
\begingroup\small
\begin{tabularx}{\linewidth}{@{}Yp{3.0cm}cp{2.4cm}cc@{}}
\toprule
Issue domain & Moderator (family) & $\beta_{\text{std}}$ & 95\% CI & $p$ & $q$ \\
\midrule
Information Integrity and Platforms & Openness (Big Five) & $+0.101$ & $[+0.041, +0.161]$ & .001 & .005 \\
Democratic Resilience and Institutions & Political trust (Pol.\ psych.) & $-0.083$ & $[-0.137, -0.029]$ & .002 & .020 \\
Foreign Policy and Geopolitics & GAL-TAN, traditional (Ideology) & $-0.075$ & $[-0.130, -0.020]$ & .008 & .030 \\
Foreign Policy and Geopolitics & Agreeableness (Big Five) & $+0.068$ & $[+0.019, +0.116]$ & .006 & .031 \\
Supranational and Regional Integration & Age (Demographics) & $-0.067$ & $[-0.117, -0.017]$ & .009 & .009 \\
Supranational and Regional Integration & Economic left-right (Ideology) & $+0.062$ & $[+0.014, +0.110]$ & .012 & .047 \\
\bottomrule
\end{tabularx}
\endgroup
\apanote{Standardised within-domain moderation slopes from ordinary least squares with HC3 robust
standard errors ($n \approx 1{,}429$ scenarios per domain), Benjamini-Hochberg corrected within each
domain-by-family cell. Only slopes surviving $q<.05$ are listed; positive values indicate greater
movability.}
\end{table}
